\documentclass{article}

\usepackage{amsmath,amssymb}
\usepackage{xurl}
\usepackage[hidelinks]{hyperref}
\setlength{\emergencystretch}{2em}

% Shared commands for notes in docs/paper-gaps/.

\DeclareUnicodeCharacter{2080}{\ifmmode{}_{0}\else\textsubscript{0}\fi}
\DeclareUnicodeCharacter{2081}{\ifmmode{}_{1}\else\textsubscript{1}\fi}
\DeclareUnicodeCharacter{2082}{\ifmmode{}_{2}\else\textsubscript{2}\fi}
\DeclareUnicodeCharacter{2099}{\ifmmode{}_{n}\else\textsubscript{n}\fi}

\newcommand{\Lean}{\textsc{Lean}}
\ExplSyntaxOn
\NewDocumentCommand{\leanid}{m}{
  \ifmmode
    \textsf{\detokenize{#1}}
  \else
    \group_begin:
    \urlstyle{sf}
    \tl_set:Nn \l_tmpa_tl {#1}
    \tl_replace_all:Nnn \l_tmpa_tl {\_} {_}
    \exp_args:NV \path \l_tmpa_tl
    \group_end:
  \fi
}
\ExplSyntaxOff
\providecommand{\tr}{\operatorname{tr}}
\newcommand{\tnleanIssueBase}{https://github.com/LionSR/TNLean/issues/}
\newcommand{\tnleanPrBase}{https://github.com/LionSR/TNLean/pull/}
\newcommand{\tnleanDocBase}{https://sirui-lu.com/TNLean/docs/}
\newcommand{\ghissue}[1]{\href{\tnleanIssueBase#1}{GitHub issue \##1}}
\newcommand{\ghpr}[1]{\href{\tnleanPrBase#1}{PR \##1}}
\newcommand{\doclink}[1]{\href{\tnleanDocBase#1.html}{\path{#1}}}

% Dirac notation and the identity operator, as in blueprint/src/macros/common.tex.
% \providecommand keeps note-local definitions authoritative.
\usepackage{dsfont}
\providecommand{\ket}[1]{|#1\rangle}
\providecommand{\bra}[1]{\langle #1|}
\providecommand{\braket}[2]{\langle #1 | #2 \rangle}
\providecommand{\ketbra}[2]{|#1\rangle\!\langle #2|}
\providecommand{\Id}{\mathds{1}}
\providecommand{\one}{\mathds{1}}
\providecommand{\C}{\mathbb{C}}
\providecommand{\MN}[1]{M_{#1}(\C)}
\providecommand{\MD}{M_D(\C)}

% External referencing for paper sources.  The build helper writes sanitized
% auxiliary files under build/paper-aux/ and excludes duplicated source labels.
\usepackage{xr}

% Import a generated paper auxiliary file with a caller-supplied label prefix.
% The candidate paths support builds from docs/paper-gaps/, the repository root,
% and build/paper-aux/ itself.
\newcommand{\importpaperaux}[2]{%
  \IfFileExists{../../build/paper-aux/#2.aux}{%
    \externaldocument[#1]{../../build/paper-aux/#2}%
  }{%
    \IfFileExists{build/paper-aux/#2.aux}{%
      \externaldocument[#1]{build/paper-aux/#2}%
    }{%
      \IfFileExists{#2.aux}{%
        \externaldocument[#1]{#2}%
      }{}%
    }%
  }%
}

% General external reference macros
\newcommand{\paperref}[2]{\ref{#1-#2}}
\newcommand{\papereqref}[2]{(\ref{#1-#2})}

% CPSV16 source (1606.00608 / MPDO-22-12-17-2.tex) external document and shortcuts
\importpaperaux{cpsv16-}{1606.00608/MPDO-22-12-17-2-tnlean}
\newcommand{\cpsvref}[1]{\paperref{cpsv16}{#1}}
\newcommand{\cpsveqref}[1]{\papereqref{cpsv16}{#1}}

% Verdict marker: \gapnote{<kind>}{<status>}. Kind is the policy
% classification (clarification, local-correction, scope-restriction,
% unfaithful, false-source, open-gap); status is open, wip (elimination
% underway), resolved, or historical. Machine-read by the site index;
% prints a verdict line.
\newcommand{\gapnote}[2]{\par\noindent\textbf{Verdict:} #1 (#2).\par}


\title{Scope of the Formal Periodic Decomposition Relative to PGVWC07}
\author{}
\date{2026-06-04}

\begin{document}
\maketitle

\section{Source theorem}

P\'erez-Garc\'ia, Verstraete, Wolf, and Cirac prove the periodic decomposition
theorem for translation-invariant MPS with periodic boundary conditions
\cite[Theorem~5, lines~849--858]{PerezGarcia2007MPSReps}.
The theorem starts with a TI state $\psi \in (\mathbb C^d)^{\otimes N}$ whose
canonical TI MPS representation has a single block with $D\times D$ matrices
$A_i$. Let
\[
  \mathcal E(X)=\sum_i A_i X A_i^\dagger.
\]
If $\mathcal E$ has $p$ eigenvalues of modulus $1$, then:
\begin{enumerate}
  \item if $p$ divides $N$, the state $\psi$ is a superposition of $p$
    $p$-periodic states, each with an MPS representation of bond dimension $D$;
  \item if $p$ does not divide $N$, then $\psi=0$.
\end{enumerate}
The proof uses the peripheral spectral structure of the CP map from
Fannes--Nachtergaele--Werner: there are orthogonal projectors $P_k$ cyclically
permuted by the matrices $A_i$, and inserting $\sum_k P_k$ into the trace
splits $\psi$ into the claimed $p$-periodic summands
\cite[lines~860--880]{PerezGarcia2007MPSReps}.

\section{Formal boundary}

TNLean proves the explicit fixed-length decomposition in
\leanid{MPSTensor.pgvwc07_periodic_stateVector_decomposition_of_dvd} and the
vanishing branch in
\leanid{MPSTensor.pgvwc07_stateVector_eq_zero_of_not_dvd}. Both declarations
use \leanid{MPSTensor.IsPeriodic}, which assumes left-canonical orientation,
irreducibility, positive period, and an exact $p$-th-root peripheral spectrum.
The printed theorem instead starts from its one-block unital canonical
representation and counts the peripheral eigenvalues of the associated
transfer map.

The periodic branch is stated for positive ring lengths because the current
coefficient decomposition is proved for a nonempty chain. The formal results
therefore give a normalized theorem surface and the complete fixed-positive-
length calculation, but not yet a literal theorem from the printed
hypotheses.

\section{Elimination plan}

Prove a source-facing equivalence from the paper's one-block unital canonical
hypotheses and peripheral-eigenvalue count to
\leanid{MPSTensor.IsPeriodic}, accounting for the adjoint normalization
orientation. Then compose this equivalence with the two fixed-length
statements above. If length zero is included in the formal source interface,
add its direct empty-chain calculation separately. No new cyclic-projector or
positive-length state-vector calculation is required.

\section{Verdict}

\textbf{Verdict: scope restriction.} The formal declarations prove the
periodic decomposition and vanishing dichotomy under the maintained
\leanid{MPSTensor.IsPeriodic} hypotheses, with positive length required in the
decomposition branch. The wrapper from the printed one-block unital
hypotheses and peripheral-eigenvalue count remains open, so the current
statements are normalized specializations of PGVWC07 Theorem~5.

\bibliographystyle{alpha}
\bibliography{references}

\end{document}